\chapter{信息访问、KL 粗化与可识别性}
\label{inf:chapter}

共同 mirror 子问题确定了要优化的分布，inner 损失确定了拟合它的几何。能否计算这个损失还取决于查询接口。完整 teacher logits、teacher Top-$J$、桶质量和实际动作的 teacher log-prob 提供不同的信息。本章先定义这些信息，再证明它们能够确定的量和仍然缺失的量。下一章讨论如何用额外查询和随机积分补回缺失贡献。

全书的 anchor、student、teacher 仍分别记为 $q,p,t$。本章分析一般 KL 时使用 $P=P_\theta$ 与冻结参考 $Q$；$P,Q$ 是一般分布的占位符，不另定义训练策略。定义 $\ell_v=\ell_v^{P,Q}=\log(P(v)/Q(v))$ 和 $\psi_v=\nabla_\theta\log P(v)$。prefix 仍为 $s$，prompt 仍为 $x$。通用 $\ell$ 与冻结 OPD reward $r=\log(t/q)$ 始终分别使用。

除专门讨论支持集边界外，本章假设有限词表上 $P(v),Q(v)>0$。记 $K=\KL(P\Vert Q)$、$K^{\mathrm F}=\KL(Q\Vert P)$，对应参数梯度为 $g=\nabla K$ 与 $g^{\mathrm F}=\nabla K^{\mathrm F}$。$\Var$ 对向量表示协方差迹。分区及所有参考查询值在求导时冻结。

\section{查询接口及其确定的信息}
\label{inf:access}

固定 prefix $s$。完整 logits 接口给出整个 $t(\cdot\mid s)$，因而所有有限求和都能精确计算。Top-$J$ 接口给出冻结支持 $S$ 及每个 $v\in S$ 的 $t(v\mid s)$；由归一化得到尾部总质量，却通常不能恢复每个尾部概率。实际动作查询给出已采样动作 $A$ 的 $t(A\mid s)$，它是一个新的坐标观测。辅助动作查询能够增加这种坐标观测；辅助动作本身按哪个分布抽取，必须另外记录。

\begin{definition}[与观测一致的参考分布]
\label{inf:compatible}
设词表分区为 $\mathcal C$，已知各桶质量 $Q_C>0$，$\sum_CQ_C=1$。仅使用桶质量时，信息类为
\begin{equation}
\mathfrak Q_{\mathcal C}=\left\{\widetilde Q:\widetilde Q(v)>0,\quad \sum_{v\in C}\widetilde Q(v)=Q_C\text{ 对每个 }C\right\}.
\end{equation}
若还观测到 token 集合 $O$ 上的概率，加入约束 $\widetilde Q(v)=Q(v)$，$v\in O$。若 $S$ 是真实 teacher Top-$J$，还加入排名约束
\begin{equation}
\widetilde Q(v)\leq\min_{u\in S}Q(u),\qquad v\notin S.
\end{equation}
相同概率的 tie 按固定规则处理。一个量可由接口确定，当且仅当它在相应信息类中取相同值。
\end{definition}

信息类只描述可观测约束。它与实际抽样分布 $\beta$、积分权重以及 student 参数化是不同对象。尾部动作的 sampled teacher log-prob 必须计入 $O$；将它忽略会人为扩大信息类。将它当成完整尾部则会人为缩小信息类。

四种 inner 所需的确定积分分别为
\begin{center}
\begin{tabular}{lll}
\toprule
inner 方法 & 所需对象 & 冻结或重算\\
\midrule
PMD-Part & $\log Z_f=\log\E_q[e^{\eta f}]$ & 整轮冻结\\
PMD-Mean & $\E_q[f]$ & 整轮冻结\\
BPO，center 表示 & $\E_q[f],\E_q[U],\E_q[\psi]$ & 信号冻结，student 项重算\\
BPO，pairwise 表示 & $U_a-U_b, f_a-f_b$ & 两个 action 上的值\\
MDPO & $\E_p[(U/\eta-f)\psi]$ & student 测度随 inner 改变\\
\bottomrule
\end{tabular}
\end{center}
这里 $U=\log(p/q)$，$\E_\mu[h]$ 是条件均值，中心化算子为 $C_\mu h=h-\E_\mu[h]$。BPO 两种表示具有相同总体目标；它无需强制先构造两个 KL center。OPD 取 $f=r$ 后，teacher 条件均值为 $\E_q[r]=-\KL(q\Vert t)$，student 条件均值为 $\E_q[U]=-\KL(q\Vert p)$。

完整 logits 消除条件 token 求和误差，仍保留 prefix 覆盖和参数化表示误差。稀疏接口是否提供精确中心，由信息类而非 estimator 名称确定。

\section{同一分区上的 KL 粗化与梯度残差}
\label{inf:sec:kl-partition-geometry}
\label{inf:sec:kl-partition}

头部精确求和、尾部合并和尾部残差抽样，都可以从同一个词表分区推导。分区同时规定可精确计算的粗粒度目标和需要抽样的细粒度残差。下面始终固定 prefix、当前参数和参考分布；只在明确写出梯度时对 student 参数求导。

\begin{definition}[词表分区与桶内分布]
令 $\mathcal C$ 是有限词表 $\mathcal V$ 的一个分区，各个单元 $C\in\mathcal C$ 非空且两两不交，并满足 $\bigcup_{C\in\mathcal C}C=\mathcal V$。分区在求导时固定。定义
\begin{equation*}
 P_C=\sum_{v\in C}P(v),\qquad
 Q_C=\sum_{v\in C}Q(v),\qquad
 \ell_C=\log\frac{P_C}{Q_C},
\end{equation*}
以及桶内条件分布和平均 score
\begin{equation*}
 P^C(v)=\frac{P(v)}{P_C},\qquad
 Q^C(v)=\frac{Q(v)}{Q_C},\qquad
 \bar \psi_C=\sum_{v\in C}P^C(v)\psi_v=\nabla\log P_C.
\end{equation*}
用 $C(v)$ 表示包含 $v$ 的唯一单元。粗粒度$\KL(P\Vert Q)$ 与桶内形状 KL 分别定义为
\begin{equation*}
 K_{\mathcal C}=\sum_{C\in\mathcal C}P_C\ell_C,
 \qquad D_C=\KL(P^C\Vert Q^C).
\end{equation*}
这里 $K_{\mathcal C}$ 是桶标签分布之间的 KL；$D_C$ 描述保留桶标签以后仍未解释的条件分布差异。
\end{definition}

\begin{theorem}[同一分区上的双向 KL 链式分解]
\label{inf:thm:kl-partition-chain}
在本章正支持假设下，$\KL(P\Vert Q)$ 和$\KL(Q\Vert P)$ 分别满足
\begin{align}
 K&=K_{\mathcal C}+\sum_{C\in\mathcal C}P_CD_C,
 \label{inf:eq:kl-partition-chain-reverse}\\
 K^{\mathrm F}&=K^{\mathrm F}_{\mathcal C}
       +\sum_{C\in\mathcal C}Q_CD^{\mathrm F}_C,
 \label{inf:eq:kl-partition-chain-forward}
\end{align}
其中
\begin{equation*}
 K^{\mathrm F}_{\mathcal C}=-\sum_CQ_C\ell_C,
 \qquad D_C^{\mathrm F}=\KL(Q^C\Vert P^C).
\end{equation*}
因此两种方向的粗粒度 KL 都是对应完整 KL 的下界。等号成立当且仅当每个单元内的条件分布相同。
\end{theorem}

\begin{proof}
对任意 $v\in C$，完整 log-ratio 可以逐点写为
\begin{equation*}
 \ell_v=\log\frac{P(v)}{Q(v)}
     =\ell_C+\log\frac{P^C(v)}{Q^C(v)}.
\end{equation*}
先乘以 $P(v)=P_CP^C(v)$，在单元内求和，再对单元求和，得到
\begin{equation*}
 \sum_vP(v)\ell_v
 =\sum_CP_C\ell_C
  +\sum_CP_C\sum_{v\in C}P^C(v)
                        \log\frac{P^C(v)}{Q^C(v)}.
\end{equation*}
这就是反向分解。将上面的逐点恒等式取负并乘以 $Q(v)=Q_CQ^C(v)$，按相同顺序求和，得到正向分解。所有质量严格为正，且每个条件 KL 非负；条件 KL 等于零恰好对应桶内条件分布相同，故得到下界及等号条件。
\end{proof}

\begin{proposition}[分区细化的标量单调性]
\label{inf:prop:kl-partition-refinement}
设 $\mathcal C'$ 细化 $\mathcal C$。记 $\mathcal C'|_C$ 为包含在 $C$ 内的细单元集合，则
\begin{equation*}
 K_{\mathcal C'}-K_{\mathcal C}
 =\sum_{C\in\mathcal C}P_C
   \KL\left(
       \left(\frac{P_{C'}}{P_C}\right)_{C'\in\mathcal C'|_C}
       \middle\Vert
       \left(\frac{Q_{C'}}{Q_C}\right)_{C'\in\mathcal C'|_C}
      \right)\geq0.
\end{equation*}
正向版本交换两组质量，并将外层权重改为 $Q_C$。因此细化使粗粒度 KL 的标量偏差绝对值和平方偏差单调下降；全部单元都是 singleton 时偏差为零。
\end{proposition}

\begin{proof}
对每个 $C$，利用 $\sum_{C'\subseteq C}P_{C'}=P_C$，将新增贡献写成
\begin{align*}
 \sum_{C'\subseteq C}P_{C'}\log\frac{P_{C'}}{Q_{C'}}
    -P_C\log\frac{P_C}{Q_C}
 &=P_C\sum_{C'\subseteq C}\frac{P_{C'}}{P_C}
       \log\frac{P_{C'}/P_C}{Q_{C'}/Q_C}.
\end{align*}
右边是 $P_C$ 乘以细单元标签之间的条件 KL。对 $C$ 求和即得所述等式。结合粗粒度下界，误差 $K-K_{\mathcal C}$ 非负且随细化下降，其平方也下降。
\end{proof}

\begin{theorem}[粗粒度 KL 的真实梯度及其残差]
\label{inf:thm:kl-partition-gradient}
记 $g_{\mathcal C}=\nabla K_{\mathcal C}$。则
\begin{align}
 g_{\mathcal C}
 &=\sum_CP_C\ell_C\bar \psi_C
  =\sum_vP(v)\ell_{C(v)}\psi_v,
 \label{inf:eq:kl-partition-gradient}\\
 g-g_{\mathcal C}
 &=\sum_CP_C\E_{P^C}[(\ell_A-\ell_C)\psi_A]
 \label{inf:eq:kl-partition-gradient-residual}\\
 &=\sum_CP_C\left\{
       \operatorname{Cov}_{P^C}(\ell_A-\ell_C,\psi_A)
       +D_C\bar \psi_C\right\}.
 \label{inf:eq:kl-partition-gradient-covariance}
\end{align}
标量与向量的协方差定义为
\begin{equation*}
 \operatorname{Cov}_{P^C}(u_A,\psi_A)
 =\E_{P^C}[(u_A-\E_{P^C}u_A)(\psi_A-\bar \psi_C)].
\end{equation*}
\end{theorem}

\begin{proof}
由于 $Q$ 固定，$\nabla Q_C=0$，而 $\nabla P_C=P_C\bar \psi_C$。逐项求导得到
\begin{align*}
 \nabla K_{\mathcal C}
 &=\sum_C\{\ell_C\nabla P_C+P_C\nabla\ell_C\}\\
 &=\sum_C\ell_C\nabla P_C+\sum_C\nabla P_C\\
 &=\sum_CP_C\ell_C\bar \psi_C.
\end{align*}
最后一步用了 $\sum_C\nabla P_C=\nabla1=0$。将 $P_C\bar \psi_C$ 展开为桶内加权 score，得到第二种表示。完整梯度为 $g=\sum_vP(v)\ell_v\psi_v$，与粗粒度梯度逐项相减，得到式~\eqref{inf:eq:kl-partition-gradient-residual}。又因为
\begin{equation*}
 \E_{P^C}[\ell_A-\ell_C]=D_C,
\end{equation*}
用 $\E[us]=\operatorname{Cov}(u,\psi)+\E[u]\E[\psi]$ 展开每个桶内期望，就得到式~\eqref{inf:eq:kl-partition-gradient-covariance}。
\end{proof}

式~\eqref{inf:eq:kl-partition-gradient-covariance} 把梯度误差分为桶内相关项和条件形状 KL 乘以平均 score 的项。标量粗化误差只包含 $D_C$ 的加权和，梯度误差还依赖 score 的方向与桶内相关性。因此标量下界的单调改进不能直接推出梯度偏差范数单调下降；梯度结论需要单独分析这些向量项。

对$\KL(Q\Vert P)$，冻结权重 $Q_C$ 使粗粒度梯度变成
\begin{equation*}
 g^{\mathrm F}_{\mathcal C}=-\sum_CQ_C\bar \psi_C,
 \qquad
 g^{\mathrm F}-g^{\mathrm F}_{\mathcal C}
 =-\sum_CQ_C\{\E_{Q^C}[\psi_A]-\bar \psi_C\}.
\end{equation*}
这里平均 student score $\bar \psi_C$ 的积分分布仍是 $P^C$；它不能替换成 teacher 条件均值。第一式来自 $\nabla\ell_C=\bar \psi_C$，第二式来自完整正向梯度 $-\sum_vQ(v)\psi_v$。

\section{粗粒度 KL 的信息解释}
\label{inf:identification}

\begin{theorem}[桶质量信息下的 KL 下包络]
\label{inf:lower-envelope}
设 $P(v)>0$，每个已知桶质量 $Q_C>0$。仅在定义~\ref{inf:compatible} 的桶质量信息类 $\mathfrak Q_{\mathcal C}$ 上，
\begin{equation}
\inf_{\widetilde Q\in\mathfrak Q_{\mathcal C}}\KL(P\Vert\widetilde Q)
=\sum_CP_C\log\frac{P_C}{Q_C}=K_{\mathcal C}.
\end{equation}
唯一达到下包络的参考分布为 $\widetilde Q(v)=Q_CP^C(v)$，$v\in C$。$\KL(Q\Vert P)$ 同样满足
\begin{equation}
\inf_{\widetilde Q\in\mathfrak Q_{\mathcal C}}\KL(\widetilde Q\Vert P)
=\sum_CQ_C\log\frac{Q_C}{P_C},
\end{equation}
并由同一个分布达到。
\end{theorem}
\begin{proof}
由双向链式分解，
\begin{align}
\KL(P\Vert\widetilde Q)&=K_{\mathcal C}+\sum_CP_C\KL(P^C\Vert\widetilde Q^C),\\
\KL(\widetilde Q\Vert P)&=K^{\mathrm F}_{\mathcal C}+\sum_CQ_C\KL(\widetilde Q^C\Vert P^C).
\end{align}
全部条件 KL 非负，因而两式分别不小于相应粗粒度目标。取 $\widetilde Q^C=P^C$ 时，构造出的 $\widetilde Q(v)=Q_CP^C(v)>0$ 满足全部桶质量约束，两式的条件 KL 同时为零。反之，正权重下条件 KL 和为零要求每个条件 KL 为零，故每个桶内分布必须相同，证明唯一性。
\end{proof}

头部 singleton、尾部单桶时，上式就是 Tail-aware KL。它将未知尾部形状差异设为零，得到桶质量信息能够保证的最小完整 KL。实际 Top-$J$ 信息还具有排名约束，因此这个锐下包络结论不能直接搬到全部真实稀疏 payload。

例如 $P=(0.1,0.8,0.1)$，已知 $Q_1=0.4$、$Q_2+Q_3=0.6$，且 $S=\{1\}$ 是真实 Top-1。仅使用质量约束的最小点为
\begin{equation}
\widetilde Q_2=0.6\frac{0.8}{0.9}=\frac8{15}>0.4.
\end{equation}
它违反 Top-1 排名。加入排名后，Tail-aware 仍是下界，但最小值严格提高：可兼容 $Q_2\leq0.4$，在闭边界最优点 $Q_2=0.4,Q_3=0.2$ 下，尾部条件分布为 $(2/3,1/3)$，与 $P^\tail=(8/9,1/9)$ 不同，条件 KL 严格为正。额外 sampled-action 概率也会收紧信息类，处理方式相同。

\begin{proposition}[相同稀疏观测不控制完整$\KL(P\Vert Q)$ 的上界]
\label{inf:unbounded}
存在相同 Top-1 payload 的正支持参考分布族，使完整$\KL(P\Vert Q)$ 无上界，而 Tail-aware 不变。
\end{proposition}
\begin{proof}
取 $P=(0.4,0.2,0.4)$ 和 $Q_\varepsilon=(0.7,\varepsilon,0.3-\varepsilon)$，$0<\varepsilon<0.15$。全部分布的 Top-1 均为 token 1，其概率 $0.7$、尾部总质量 $0.3$ 相同。可是
\begin{equation}
\KL(P\Vert Q_\varepsilon)
=0.4\log\frac{0.4}{0.7}+0.2\log\frac{0.2}{\varepsilon}
+0.4\log\frac{0.4}{0.3-\varepsilon}\longrightarrow+\infty.
\end{equation}
第一项固定，第三项趋向有限值，第二项发散。此例需要至少两个仍未知的尾部坐标。若 payload 已查询其中一个，尾总质量会确定另一个，信息类不再保留这个自由度。
\end{proof}

\section{非线性积分与无偏性陈述的边界}
\label{inf:nonlinear}

\begin{proposition}[正的 plug-in partition 估计的 Jensen 偏差]
\label{inf:jensen}
若 $\widehat Z>0$、$\E\widehat Z=Z$ 且 $\E|\log\widehat Z|<\infty$，则
\begin{equation}
\E[\log\widehat Z]\leq\log Z.
\end{equation}
严格凹性使等号当且仅当 $\widehat Z$ 几乎处处为常数。
\end{proposition}
\begin{proof}
在 $Z>0$ 处，$\log$ 的切线满足 $\log z\leq\log Z+(z-Z)/Z$，且 $z\ne Z$ 时严格。取期望并利用 $\E(\widehat Z-Z)=0$ 即得结论；若偏离 $Z$ 的事件具有正概率，严格不等式在该事件成立，期望不等式也严格。
\end{proof}

这个命题分析 $\log\widehat Z$，并没有证明所有 $\log Z$ 估计器都不能无偏。双 token 词表中，若 anchor 全分布已知，查询任一 token 的 teacher 概率后，另一概率由归一化确定。因此一个 sampled-token 查询就能精确计算任意 $\eta$ 的 OPD partition。

更一般，设 $|\mathcal V|=D$，$m\geq D$，从已知正支持 anchor 独立抽样 $m$ 次并查询每个动作的 teacher 概率。令 $E$ 为样本覆盖整个词表的事件，$c=\Pr_q(E)>0$。$c$ 只依赖已知 anchor 与样本数。当 $E$ 成立时可恢复 teacher 全分布。定义
\begin{equation}
\widehat L=\mathbf1_E\frac{\log Z}{c},
\end{equation}
并在 $E$ 未发生时直接返回零，不调用未知 $\log Z$。逐项计算得到
\begin{equation}
\E\widehat L=\log Z,\qquad
\Var(\widehat L)=(c^{-1}-1)(\log Z)^2.
\end{equation}
该构造通常具有不可接受的成本和方差，却证明了“所有无偏 log-partition 估计都不存在”这一陈述不成立。实际方法应分别报告信息条件、估计器形式、偏差和成本。

\section{粗化、条件期望与真正的 Rao--Blackwell 化}
\label{inf:rao-blackwell}

在 $A\sim P$ 下，若知道桶内形状 KL，真实条件期望为
\begin{equation}
\E_P[\ell_A\mid C(A)=C]=\ell_C+D_C.
\end{equation}
由全方差公式，
\begin{equation}
\Var_P(\ell_A)
=\E_P[\Var_P(\ell_A\mid C(A))]
+\Var_P(\E_P[\ell_A\mid C(A)]).
\end{equation}
因此用真实条件期望替代原始样本不会增加方差。Tail-aware 仅知道 $\ell_C$，未知的 $D_C$ 没有自动出现；HETS 将头部求和并留下尾部事件修正，也不是上述条件期望。后者的方差是否下降，要由残差二阶矩直接判断。向量情形应用全协方差公式并取迹，逻辑完全相同。

